\section{Problem Statement / Limitations}

Practical laboratory examinations are often conducted in environments where the technical challenge of the assessment is matched by an equally serious management challenge. In many academic institutions, there is no single, unified platform that coordinates scheduling, invigilation, access control, workstation-level tracking, incident recording, and post-exam reporting. As a result, critical parts of the process remain manual, fragmented, or weakly documented, making the examination workflow harder to supervise and more difficult to justify when concerns arise. One of the most significant limitations of this current environment is the limited ability of invigilators to monitor many students simultaneously with confidence. In a crowded computer lab, human attention alone is insufficient to maintain real-time awareness of who has started, who is inactive, who may be experiencing technical issues, and who may be deviating from expected behavior, which is further aggravated by the lack of a consolidated dashboard to support situational awareness \cite{mit_cbt}.

A second major limitation concerns the lack of evidence preservation and operational traceability during high-stakes practical testing. In many practical exams, suspicious activity, submission problems, or timing disputes are extremely difficult to reconstruct after the session has ended because there is no automated event logging. Without structured logs, screen focus shift timelines, and organized incident records, academic decisions and grading overrides must depend on fragmented personal observations rather than well-preserved, auditable evidence. This issue is compounded by the high administrative workload caused by utilizing separate tools for manual coordination. Scheduling, attendance tracking, active workstation supervision, file collection, and reporting are handled via isolated channels, which significantly increases operational overhead and the possibility of data inconsistency across semesters.

The proposed project itself also has specific development and operational limitations that define its scope. As a final year project (FYP), it will emphasize a well-structured prototype and an academically defensible design rather than claiming a full enterprise-scale deployment across all university departments. Advanced AI-driven anomaly detection may initially be represented through configurable monitoring rules and alert triggers rather than a production-trained deep learning predictive model. Similarly, large-scale live deployment constraints such as network latency variance, legacy institutional database integration complexity, and long-term security maintenance policies will remain outside the immediate prototype boundary. These boundaries are explicitly acknowledged to ensure that the project remains technically realistic and achievable within academic timeframes while still delivering meaningful operational value.
